\section{Segunda iteración de la Fase 2: Preparando la base}
\label{sec:PreparandoBase}

El análisis de la iteración anterior definió el sistema reducido con sus tres componentes y las reglas que los separan (un solo escritor de la lista de segmentos, banderas \emph{visible} y \emph{cambiar}, motor gráfico lector puro). Esta iteración comienza el desarrollo, que se desarrollará en \textbf{tres pasos}: primero la base que aquí se presenta, después montaremos mejro el director y el motor de juego mínimo con la secuencia fija, y por último el objetivo real de esta fase, el motor gráfico real. El desarrollo arranca por lo más simple a propósito: montar el \emph{esqueleto completo} ---director, dos motores, bucle, pausa, y el arranque desde la máquina--- con todo su contenido con respuestas fijas, para verificar que el esqueleto en sí funciona antes de nada. Es un paso deliberadamente simple, pero prefiero ir paso a paso con funcionalidad mínima, e ir ampliándola, que lanzarme a un desarrollo de varios temas a la vez, mejor avanzar poco a poco sobre algo que se ha probado.

\subsection{Qué se construye}
\label{subsec:QueSeConstruyeF2I2}

\begin{itemize}
	\item \textbf{Director funcional.} Aunque sus llamadas a otras funciones solo sean un esqueleto vacío, el director de esta iteración ya está bien definido y listo para futuras ampliaciones: mantiene un \textbf{contador de ticks} que incrementa en cada ciclo, lo \textbf{presenta en pantalla} (en la barra de estado) como referencia, llama al motor del juego, llama al motor gráfico, y se hace la \textbf{pausa final} encadenando el temporizador, y continúa en bucle. El ritmo y el orden de llamada ya son los definitivos: \texttt{tick() $\rightarrow$ estado $\rightarrow$ presentar $\rightarrow$ esperar}.
	\item \textbf{Motores sin funcionalidad.} El motor del juego y el motor gráfico existen ya como ficheros y objetos globales con su nombre y su hueco, pero su único cometido es \emph{anunciar que viven}: al recibir el control escriben un mensaje ---«motor del juego» y «motor gráfico» respectivamente--- para verificar que el director los llama en orden y una vez por tick. No tocan segmentos ni pantalla: cualquier cosa que hagan más allá de anunciarse sería anticipar pasos que pertenecen a las iteraciones posteriores.
	\item \textbf{Arranque desde la máquina.} El \texttt{index.html} ya activa el botón \textbf{Start} que lanza el director. El arranque es deliberadamente manual: no basta con cargar el cartucho, hay que accionar la máquina, como en la real. Mientras el simulador no tenga más lógica de sesión, arrancar ya en marcha automático confundiría la carga de un cartucho con el juego en sí.
\end{itemize}

\subsection{Detalles de realización}
\label{subsec:DetallesF2I2}

\begin{itemize}
	\item \textbf{Ficheros nuevos, nombres definitivos}: los tres componentes nacen con su nombre final de fichero y de objeto global (\texttt{DirectorJuego}, \texttt{juegomotorestados}, \texttt{JuegoMotorGrafico}), para que las iteraciones siguientes los llenen sin renombres.
	\item \textbf{El contador en la barra de estado}: el número de tick se muestra con el mecanismo de avisos ya existente, sin crear canales de depuración nuevos. Ver el contador avanzar uno a uno al ritmo de la pausa ES la prueba de que el reloj funciona.
	\item \textbf{Pausa encadenada}: el director repite con \texttt{setTimeout} encadenado tras cada tick, nunca \texttt{setInterval}, de modo que un tick no empieza antes de que acabe el anterior. El bucle se detiene limpio si se lanza de nuevo (no hay dos relojes simultáneos).
	\item \textbf{Orden de scripts} en \texttt{index.html}: los ficheros nuevos entran tras los existentes, \texttt{juegomotorestados.js} y \texttt{JuegoMotorGrafico.js} antes que \texttt{directorjuego.js}, manteniendo la convención del script clásico sin módulos, compatible con \texttt{file://}.
\end{itemize}

El resultado de la iteración es observable y modesto: pulsar Start y ver, al ritmo de la pausa, el contador de ticks avanzar en la barra de estado mientras la consola alterna «motor del juego» y «motor gráfico». Con el esqueleto latiendo, la siguiente iteración desarrolla el motor del juego básico, con la secuencia fija de segmentos.

\subsection{Cambios de nombres}
 El \texttt{Director} de la Fase 1 (orquestador del editor) puede confundirse con este nuevo director del juego, para evitarlo lo aclararemos usando nuevos nombres de los ficheros:

\begin{itemize}
  	\item Los elementos para la autoría de los cartuchos se renombrar añadiendo delante del nombre la palabra \textbf{Base}
  	\item Los motores del juego usarán delante la palabra \textbf{Juego}
  	\item Como la parte de creación de cartuchos y la del simulador deben ser independientes, creo un nuevo fichero para mantener en el las listas de objetos que son comunes a ambas partes, de esta manera se independizan al máximo, ya que es un fichero que se comporta como un \emph{INCLUDE} común cuando es necesario.
\end{itemize}
  
  De esta manera no hay confusión entre los usos de cada uno.
  
\subsection{El puente de datos: ComunListas}
\label{subsec:ComunListas}
  
Como los motores necesitan utilizar la lista de segmentos, pero la autoría de cartuchos y los motores del simulador no deben verse entre sí, ya que son partes independientes del proyecto, y como ambas etapas necesitan \emph{los mismos datos} del cartucho una vez cargado
  (la lista de segmentos en memoria y la imagen de fondo), lo separo creando un \textbf{fichero propio que no pertenece a ninguna de las dos etapas}, \texttt{ComunListas.js}, que pueden manejar ambas partes sin conflictos.
  
Este fichero se comporta conceptualmente como un \emph{INCLUDE} en los módulos que lo necesiten utilizar.
  
  \begin{itemize}
  	\item \textbf{Es un puente, no un motor}: solo almacena referencias y datos, no interpreta
  	nada. La autoría \emph{escribe}: al cargar un GWI, \texttt{cargar()} deja las listas
  	(filtrando los fantasmas de unión, tipo 99) y \texttt{fijarFondo()} cuelga el
  	\texttt{img\_f}; al recargar, \texttt{limpiar()} vacía todo. El juego \emph{lee}:
  	\texttt{segmentos()} (por referencia, viva) y \texttt{fondo()}.
  	
  	\item \textbf{Contrato intacto}: «la lista ES el estado» sigue siendo de los motores, 
  	 \texttt{ComunListas} solo guarda referencias.
  	
  	\item \textbf{Quién escribe y quién lee en cada etapa}: la autoría escribe (carga del
  	GWI: \texttt{BaseDescomponer.cargarListas} delega en el puente, \texttt{reset} lo
  	limpia, \texttt{BaseDirector} fija el fondo); el motor de estados escribe las banderas
  	\texttt{visible}/\texttt{cambiar} en los objetos de la lista del puente; el motor
  	gráfico lee fondo, posiciones e imágenes y \emph{nunca} escribe.
  \end{itemize}
  
De esta manera independizamos ambas partes y evitamos duplicidad, un fichero común en lugar de copiar el código. Solo hay que tener cuidado de que si en una fase se modifica, en las anteriores debe seguir sirviendo, o estas se deben modificar de manera adecuada.
  
  \begin{itemize}
  	\item \textbf{\texttt{ComunListas.js} (nuevo)}: API \texttt{cargar}, \texttt{limpiar},
  	\texttt{fijarFondo}, \texttt{segmentos}, \texttt{sprites}, \texttt{fondo}. Script
  	clásico, compatible con \texttt{file://}, cargado en primer lugar en
  	\texttt{index.html} (tras \texttt{jszip.min.js}).
  	 	
  	\item \textbf{\texttt{index.html}}: debe ubicar \texttt{<script src="designer/ComunListas.js">} como primer script del proyecto, común a todos.

  	\item \textbf{\texttt{BaseDirector.js}}: junto a
	\texttt{BaseGraficos.fijarFondo(imgF)}, fija también el fondo en el puente
	(\texttt{ComunListas.fijarFondo}).

  	\item \textbf{\texttt{BaseDescomponer.js}}: \texttt{cargarListas} y
  	\texttt{segmentosEnMemoria} pasan a delegar en el puente, en lugar de usarlos como variables
  	propias de listas); \texttt{reset} limpia el puente. Las funciones de pintado del
  	aparato dejan de exponerse para el motor (el motor ya no las usa).
  	
 	\item \textbf{\texttt{JuegoMotorEstados.js}}: El futuro motor de estados del juego 
 	 puede leer y cambiar la información de los segmentos que existen en \texttt{ComunListas}. 
  	
  	\item \textbf{\texttt{JuegoMotorGrafico.js}}: El futuro motor gráfico leerá la lista desde \texttt{ComunListas}, sin alterarlas. 
  	
  \end{itemize}
  
  \clearpage
  
  \subsection{Estructura de Directorio}
  \label{subsec:EstructuraDirectorioF2I1}
   Con estos cambios, esta es la estructura de directorio actual:
  
 
 \dirtree{%
 	.1 GW/.
 	.2 index.html.
 	.2 designer/.
 	.3 jszip.min.js.
 	.3 estilos.css.
 	.3 ComunListas.
 	.3 BaseDirector.js.
 	.3 BaseDescomponer.js.
 	.3 BaseExtractor.js.
 	.3 BaseGraficos.js.
 	.3 JuegoDirector.js.
 	.3 JuegoMotorEstados.js.
 	.3 JuegoMotorGrafico.js.
 	.2 cartucho/.
 	.3 parchute.json.
 	.3 parachuite\_f.png.
 	.3 parachuite\_m.png.
 	.3 parachuite\_s.png.
 	.2 parachute\_gwi.zip.
  	.2 setup2.json.
 	.2 setup4.json.
 	.2 setup5.json.
 }
 
 \subsection{Ejecución}
 
 Cargamos un setup, luego un cartucho, y al pulsar sobre START vemos que el contador de ticks aumenta y aparecen los mensajes en la consola.
 
 \begin{figure}[H]¢
 	\centering
 	\includegraphics[width=1\linewidth]{imagenes/propias/C05-I019}
 	\caption[]{Ejecución del motor del juego básico.}
 	\label{fig:c05-i019}
 \end{figure}
 
\clearpage